\originalchapter{习题课一: 数学工具}
本章题源为\href{https://ocw.mit.edu/courses/8-02-physics-ii-electricity-and-magnetism-spring-2007/resources/solving01/}{Problem Solving 1}, PDF第1--12页. 原件包含讲解、空白答题页和重复题干, 下列将同题合并, 保留P1--P7全部小问. 所有解答为AI独立解答. 原件指向的Shockwave坐标交互只作历史引用, 本册用静态定义完成所需观察.

\begin{exercise}\label{ps1:1}
P1. 柱坐标和球坐标的单位坐标轴, 与直角坐标轴有什么根本不同?
\end{exercise}
\begin{solution}
直角单位矢量在不同位置方向不变; 柱坐标径向、方位向以及球坐标径向、极角向、方位向的单位矢量随位置改变方向. 例如$\uv s=(\cos\phi,\sin\phi,0)$随$\phi$转动, 不是整空间一根固定的$x$轴. 柱坐标轴上方位角不唯一, 球心径向不唯一, 这些是坐标奇点而非物理场必然发散.
\end{solution}

\begin{exercise}
P2. 对$V(x,y,z)=\kc q/\sqrt{x^2+y^2+z^2}$, 分别求(a)$\partial_xV$, (b)$\partial_yV$, (c)$\partial_zV$.
\end{exercise}
\begin{solution}
设$r=\sqrt{x^2+y^2+z^2}>0$. 链式法则给$\partial_x(r^{-1})=-r^{-2}x/r=-x/r^3$. 因而三答案依次为
\[
\partial_xV=-\frac{\kc qx}{r^3},\qquad
\partial_yV=-\frac{\kc qy}{r^3},\qquad
\partial_zV=-\frac{\kc qz}{r^3}.
\]
\end{solution}

\begin{exercise}
P3. 对同一$V$, (a)求梯度; (b)用$r$和$\uv r$改写; (c)说明梯度是否总指向$V$的“上坡”方向.
\end{exercise}
\begin{solution}
(a)$\nabla V=-\kc q(x\uv x+y\uv y+z\uv z)/r^3$; (b)$\nabla V=-\kc q\uv r/r^2$; (c)在梯度非零处, 方向导数$\nabla V\cdot\uv u$在$\uv u$与梯度同向时最大, 所以是上坡方向. 对正$q$向内, 负$q$向外. $q=0$时没有唯一上坡方向, 原点不在公式定义域.
\end{solution}

\begin{exercise}
P4. 半径$R$、高$H$的圆柱体, 用(a)三重小块积分, (b)以半径$r$、厚$\dd r$、高$\dd z$的环面积积分, (c)以高$H$的柱壳积分, (d)以厚$\dd z$的圆盘积分, 分别写出体积.
\end{exercise}
\begin{solution}
四种写法分别为
\begin{align*}
\mathcal V&=\int_0^H\int_0^{2\pi}\int_0^Rr\dd r\dd\phi\dd z
=\int_0^H\int_0^R2\pi r\dd r\dd z\\
&=\int_0^R2\pi rH\dd r=\int_0^H\pi R^2\dd z=\pi R^2H.
\end{align*}
后面三式只是把已知均匀的角度或高度积分先算完, 没有改变被加总的体积.
\end{solution}

\begin{exercise}
P5. 求总电荷: (a)半径$R$的圆环, 均匀线密度$\lambda$; (b)半径$R$的实心球, 体密度$\rho(r)=\rho_Rr/R$.
\end{exercise}
\begin{solution}
(a)$Q=\lambda\int\dd\ell=2\pi R\lambda$. (b)每个球壳电荷$\dd q=(\rho_Rr/R)4\pi r^2\dd r$, 故$Q=(4\pi\rho_R/R)\int_0^Rr^3\dd r=\pi\rho_RR^3$. $\rho_R$是球表面的密度, 不是平均密度.
\end{solution}

\begin{exercise}
P6. 电荷$q$在匀强场$\ee=E_0\uv y$中, 沿原点为中心、半径$R$的$xy$平面半圆, 从$(0,R)$移到$(0,-R)$. 电场作功多少?
\end{exercise}
\begin{solution}
$W=qE_0\int\dd y=qE_0(-R-R)=-2qE_0R$. 也可用半圆参数积分, 但匀场的功只依赖端点的$y$差, 半圆选左侧或右侧均不改变答案.
\end{solution}

\begin{exercise}
P7. (a)均匀电场$\ee=a\uv x+b\uv y$穿过面积$A$的平面, 求(i)法向沿正$y$的$xz$平面通量, (ii)法向沿正$z$的$xy$平面通量. (b)半径$R$、高$h$、轴沿$z$的圆柱, 在$\ee=E_0\uv y$中, 求$y>0$一半侧面的向外通量. 写法向和面积元.
\end{exercise}
\begin{solution}
(a)(i)$\Phi=bA$; (ii)$\Phi=0$. (b)设$0<\phi<\pi$, $\uv n=\cos\phi\uv x+\sin\phi\uv y$, $\dd A=R\dd\phi\dd z$. 所以$\Phi=E_0R\int_0^h\dd z\int_0^\pi\sin\phi\dd\phi=2E_0Rh$. 它等于场强乘侧面对垂直场方向的投影面积$2Rh$, 不是乘半圆柱侧面的实际面积$\pi Rh$.
\end{solution}
